% LaTeX preamble for the OSPA spec PDF, following the 2024 OCP Base
% Specification Template (v1.6). \ospatitle and \ospadate are set from the
% document by ocp-spec.lua.
\usepackage{xcolor}
\definecolor{ocpgray}{HTML}{434343}
\definecolor{ocpblue}{HTML}{1F4E79} % hyperlinks
\providecommand{\ospatitle}{}
\providecommand{\ospadate}{}

% Template margins: 1" left/right/bottom; body text starts 1.44" from the top,
% with the header 0.35" and the footer 0.5" from the page edge.
\usepackage{geometry}
\geometry{top=1.44in, bottom=1in, left=1in, right=1in,
  headheight=14pt, headsep=0.9in, footskip=0.5in}

% Template header: centered "Open Compute Project • <title>"; footer: date on
% the left, "Page N" on the right; 10pt, no rules. The title page has neither.
\usepackage{fancyhdr}
\pagestyle{fancy}
\fancyhf{}
\fancyhead[C]{\small Open Compute Project \textbullet{} \ospatitle}
\fancyfoot[L]{\small Date: \ospadate}
\fancyfoot[R]{\small Page \thepage}
\renewcommand{\headrulewidth}{0pt}
\renewcommand{\footrulewidth}{0pt}
% Chapter-opening pages (e.g. the ToC) use the same header/footer.
\fancypagestyle{plain}{}

% Headings carry their numbers in the Markdown text. LaTeX numbering must
% still be enabled: hyperref only creates a link target for a heading when
% its section counter steps, so with numbering off the ToC entries and PDF
% bookmarks point at whatever anchor preceded the heading. Enable the
% counters (after pandoc's own secnumdepth setting) and hide the generated
% numbers in both the headings and the ToC.
\AtBeginDocument{\setcounter{secnumdepth}{3}}
\setcounter{tocdepth}{3}
\AtBeginDocument{\renewcommand{\numberline}[1]{}}
% Headings use LaTeX's standard \@startsection rather than titlesec. The link
% target is placed inside the heading line (\@seccntformat is typeset at the
% start of it, in place of the hidden number): a target placed before the
% heading can be separated from it by a page break, which linked headings
% that start a new page to the bottom of the previous page.
% Template heading styles: Heading 1 14pt bold, Heading 2 12pt bold,
% Heading 3 14pt regular in gray #434343; spacing before/after 20/6pt,
% 18/6pt and 16/4pt.
\makeatletter
\renewcommand{\@seccntformat}[1]{\protect\phantomsection}
\renewcommand{\section}{\@startsection{section}{1}{\z@}%
  {-20pt \@plus -4pt \@minus -2pt}{6pt}%
  {\normalfont\fontsize{14}{17}\selectfont\bfseries}}
\renewcommand{\subsection}{\@startsection{subsection}{2}{\z@}%
  {-18pt \@plus -4pt \@minus -2pt}{6pt}%
  {\normalfont\fontsize{12}{14.5}\selectfont\bfseries}}
\renewcommand{\subsubsection}{\@startsection{subsubsection}{3}{\z@}%
  {-16pt \@plus -4pt \@minus -2pt}{4pt}%
  {\normalfont\fontsize{14}{17}\selectfont\color{ocpgray}}}
\makeatother

% Figures stay where they appear in the text, with centered "Figure N: ..."
% captions below them, as in the template.
\usepackage{float}
\floatplacement{figure}{H}
\usepackage[labelsep=colon, justification=centering]{caption}

% Tables: ocp-tables.py gives every table a full grid with a bold header row,
% as in the template, which uses 1pt black borders. Add a little padding so
% cell text does not touch the rules.
\setlength{\arrayrulewidth}{1pt}
\setlength{\extrarowheight}{2pt}

% Keep wide images and tables inside the margins.
\setlength{\emergencystretch}{3em}
